% !TeX root = ../教材题库_全部条目_转写工作稿.tex
% Source: XJ-B1-C03-LAB-P073-Y0145
% 第3章 / 3.1 函数 / 教材 p.73 / PDF 第80页
% OCR未校对：公式与文字均可能有误；未完成数学TeX转写；禁止直接发给学生
% 原题及相邻原解截图（核对用）：../题目原图/XJ-B1-C03-LAB-P073-Y0145-01.webp
% 本文件只在转写完成后作为题干使用；OCR 条目尚不是可靠题干。
\begin{bankitem}{用计算机作函数图象和列函数表}{第3章 / 3.1 函数 / 教材 p.73 / PDF 第80页}{XJ-B1-C03-LAB-P073-Y0145}{OCR未校对：公式与文字均可能有误；未完成数学TeX转写；禁止直接发给学生}
\begin{ocrdraft}
[教材p.73 / PDF第80页]
\par
数学实验
\par
用计算机作函数图象和列函数表
\par
函数的图象能够帮助我们了解函数的性质．我们可以使用计算机技术，根据函
\par
数解析式和数据作出函数的图象．使用“超级画板”或其他具有类似功能的软件，
\par
在计算机上作函数的图象既快捷又准确．
\par
在右键菜单里单击“函数或参数方程曲线”，在出现的对话框里（如图1）输入函
\par
数表达式“（1/2）”，参数范围取0到10，点数取101，单击“确定”，函数3，一，在
\par
的图象就出现了（如图2）．
\par
对象餉落
\par
函庫笔充置《文事|
\par
0=f回
\par
0方程
\par
0和坐标=《@〕
\par
间点小值伸应生标〕印一瞓铲点数\ensuremath{\cdot}LiOi\ocrglyph{00AF}囗折段
\par
囗画囗点的大小，一囗斷段
\par
定《取消
\par
（丿
\par
用鼠标在曲线附近单击左键，曲线变色，表示计算机知道你选择了它．
\par
执行菜
\par
单命令“插入》表格”，屏幕上就出现了对应的函数表．
\par
3，
\par
3，
\par
0．00000．00000．99991.00002，000011423，00001.73204，00002．00005，00002，23606，00002，44947以〕000
\par
0，l(〕000，31621.10001.04882，10001.M913，l(〕001.76064，10002，02485，l(〕00巳25836，10002，46987．l(〕00
\par
0．20000．11721.20001.09542，2000118323，20001.78884，20002，04935．20002，28036，20002，48997，2000
\par
0．30000．5771.30001.14012，30001.51653，30001.81654，30002，07365，30002，30216，30002，50997，3000
\par
0．40000．63241.40001.18322，0001.54913，4000184394．40002．09765，40002，32376，40002，52987，4000
\par
0．50000．70711.50001.22472，50001.58H3，50001.87084，50002，12135，5000巳引巧26，50002，54957，5000
\par
0．60000．77451.60001.26492，60001.61243，60001.89734，60002，14475，60002，36646．60002．56907，6000
\par
0，70000，83661.70001.30382，70001.6313，70001.92354，70002，16795，70002，38746．70002．58847，7000
\par
0．80000．89441.80001.34162，80001.67333，80001.94934，80002，19085，80002，10836，80002．60767，8000
\par
0．90000．94861.9000137842，90001.70293．90001.97484．90002．21355，90002，42896．90002．62677，9000
\par
这时表上的格子不够，可以用鼠标指着表格按右键，在右键菜单中单击“属
\par
性”，打开表格的属性设置对话框，在对话框下部把表格的行数改为11，列数改为
\par
20，单击“确定”即可．
\end{ocrdraft}
\end{bankitem}
